% from main.tex
%%%%%%%%%%  scan idx 111  |  folio 98  %%%%%%%%%%
\textbf{Définition 2.3.} \emph{Sous les hypothèses de (2.2), nous dirons que \(F\in\Ob\Dpx(\mathcal C_S)\) est} \(n\)-\(\mathcal C_0\)-pseudo-cohérent \emph{si \(F\) vérifie les conditions équivalentes b) de (2.2). Nous dirons que \(F\) est} \(\mathcal C_0\)-pseudo-cohérent \emph{si \(F\) est \(n\)-\(\mathcal C_0\)-pseudo-cohérent pour tout \(n\).}

Quand cela ne pourra prêter à confusion, nous nous permettrons parfois de dire \(n\)-pseudo-cohérent, pseudo-cohérent, au lieu de \(n\)-\(\mathcal C_0\)-pseudo-cohérent, \(\mathcal C_0\)-pseudo-cohérent.

\textbf{Remarque 2.3.1.} Soit \(F\in\Ob\Dpx(\mathcal C_S)\) un complexe pseudo-cohérent. L'ensemble des \(X\) au-dessus de \(S\) tels qu'il existe un quasi-isomorphisme \(E\longrightarrow F_X\) avec \(E\) strictement pseudo-cohérent (2.1) n'est pas en général un raffinement de \(S\). Toutefois, nous verrons plus tard (2.7. a) et Exp.II) des cas importants où il en est ainsi.

La démonstration de (2.2 a)) donne le résultat suivant :

\textbf{Proposition 2.4.} \emph{Soient \(n,i\in\mathbb Z\), et soit \(u:E\to F\) un \(i\)-quasi-isomorphisme de \(\Cpx(\mathcal C_S)\), avec \(F\) \(n\)-pseudo-cohérent et \(E\) strictement \(n\)-pseudo-cohérent. Alors, quitte à raffiner \(S\), il existe un triangle commutatif de \(\Cpx(\mathcal C_S)\)}
\[
\begin{tikzcd}
E \arrow[r,"v"] \arrow[d,"u"'] & E' \arrow[dl,"u'"]\\
F & ,
\end{tikzcd}
\]
\emph{où \(E'\) est strictement \(n\)-pseudo-cohérent, \(u'\) un quasi-isomorphisme et \(v\) un monomorphisme induisant l'identité en degré \(\geq i\).}

\textbf{Proposition 2.5.} a) \emph{Si \(F\in\Ob\Dpx(\mathcal C_S)\) est \(n\)-pseudo-cohérent (resp. pseudo-cohérent), \(F[m]\) est \((n-m)\)-pseudo-cohérent (resp. pseudo-cohérent).}

b) \emph{Soit}
